\documentclass[12pt,a4paper]{article}
\usepackage[spanish,es-nodecimaldot]{babel}
\usepackage[utf8]{inputenc}
\usepackage[T1]{fontenc}
\usepackage{geometry}
\usepackage{hyperref}
\usepackage{listings}
\usepackage{xcolor}
\geometry{margin=2.5cm}
\definecolor{codegray}{gray}{0.95}
\lstset{language=R,basicstyle=\ttfamily\small,backgroundcolor=\color{codegray},frame=single,breaklines=true}
\title{Modelos híbridos}
\author{Unidad 1}
\date{2026}
\begin{document}
\maketitle
\section{Definición}
Un modelo híbrido combina dos o más técnicas para aprovechar sus fortalezas. Puede integrar selección de variables con un clasificador, varios predictores en un ensemble o métodos estadísticos con aprendizaje automático.
\section{Diseño}
Un modelo híbrido debe justificar qué aporta cada componente, definir cómo se combinan las salidas y compararse con modelos individuales mediante la misma partición de datos.
\section{Aplicaciones}
Se utiliza en predicción de demanda, clasificación médica, crédito, detección de intrusiones y mantenimiento predictivo.
\section{Tipos de combinación}
La combinación puede ocurrir antes del modelado, mediante selección o transformación de variables; durante el modelado, mediante ensambles; o después del modelado, mediante promedio, votación o stacking de predicciones.
\section{Ventajas y costos}
Los modelos híbridos pueden reducir la varianza y capturar relaciones complejas, pero aumentan la complejidad, el tiempo de entrenamiento y la dificultad de explicar el resultado. La combinación debe justificarse con una comparación contra cada modelo individual.
\section{Diseño experimental}
Usar la misma división de datos, fijar la semilla, evitar fuga de información y reportar métricas, costo computacional y sensibilidad a cambios en los datos.
\section{Ejemplo en R}
\begin{lstlisting}
set.seed(10)
datos <- data.frame(x1 = rnorm(60), x2 = rnorm(60))
datos$clase <- factor(ifelse(datos$x1 + datos$x2 > 0, "si", "no"))
indices <- sample(seq_len(nrow(datos)), size = 0.7 * nrow(datos))
entrenamiento <- datos[indices, ]
prueba <- datos[-indices, ]
modelo_a <- glm(clase ~ x1 + x2, data = entrenamiento, family = binomial)
modelo_b <- glm(clase ~ x1 * x2, data = entrenamiento, family = binomial)
prob_a <- predict(modelo_a, prueba, type = "response")
prob_b <- predict(modelo_b, prueba, type = "response")
prob_promedio <- (prob_a + prob_b) / 2
prediccion <- factor(ifelse(prob_promedio > 0.5, "si", "no"), levels = levels(prueba$clase))
mean(prediccion == prueba$clase)
\end{lstlisting}
\section{Referencias iniciales}
\begin{itemize}
\item Francis, B. K. y Babu, S. S. (2019). Predicting Academic Performance Using a Hybrid Data Mining Approach. DOI: \url{https://doi.org/10.1007/s10916-019-1295-4}.
\item Koutanaei, F. N. y colaboradores (2015). A hybrid data mining model for credit scoring. DOI: \url{https://doi.org/10.1016/j.jretconser.2015.07.003}.
\end{itemize}
\end{document}
